\documentclass{article}
\usepackage{geometry}
\geometry{margin=1.5cm, vmargin={0pt,1cm}}
\setlength{\topmargin}{-1cm}
\setlength{\headheight}{12.5pt}
\setlength{\paperheight}{29.7cm}
\setlength{\textheight}{25.3cm}

% useful packages.
% \usepackage[UTF8]{ctex}
\usepackage{fancyhdr}
\usepackage{layout}
\usepackage{tikz}
\usetikzlibrary{arrows.meta}

% import common macros
% % % % % % % % % % % % % % % % % % % % % % % % % % % % % % % %
% Common Macros for LaTeX
% % % % % % % % % % % % % % % % % % % % % % % % % % % % % % % %

% Useful packages
\usepackage{amsfonts}
\usepackage{amsmath}
\usepackage{amssymb}
\usepackage{amsthm}
\usepackage{enumerate}
\usepackage{graphicx}
\usepackage{multicol}
\usepackage[ruled,linesnumbered]{algorithm2e}

% Math operators and common symbols
\newcommand{\dif}{\mathrm{d}}
\newcommand{\innerProd}[1]{\left\langle #1 \right\rangle}
\newcommand{\difFrac}[2]{\frac{\dif #1}{\dif #2}}
\newcommand{\pdfFrac}[2]{\frac{\partial #1}{\partial #2}}
\newcommand{\T}{\mathrm{T}}
\newcommand{\norm}[1]{\left\| #1 \right\|}
\newcommand{\abs}[1]{\left| #1 \right|}

% Vector and Matrix shortcuts
\newcommand{\vecTwo}[2]{
  \begin{bmatrix}
    #1 \\
    #2
\end{bmatrix}}
\newcommand{\vecThree}[3]{
  \begin{bmatrix}
    #1 \\
    #2 \\
    #3
\end{bmatrix}}
\newcommand{\vecTwoH}[2]{
  \begin{bmatrix}
    #1 & #2
  \end{bmatrix}
}
\newcommand{\vecThreeH}[3]{
  \begin{bmatrix}
    #1 & #2 & #3
  \end{bmatrix}
}

% Common fields
\newcommand{\R}{\mathbb{R}}
\newcommand{\C}{\mathbb{C}}
\newcommand{\Z}{\mathbb{Z}}
\newcommand{\N}{\mathbb{N}}

% Solutions and proofs
\newenvironment{solution}{
\begin{proof}[解]}{
\end{proof}}

% Utility to get environment variables (requires lualatex)
\newcommand{\getEnv}[2]{\directlua{tex.print(os.getenv("#1") or "#2")}}


% Name and uID
\newcommand{\name}{\getEnv{NAME}{NAME}}
\newcommand{\uid}{\getEnv{UID}{UID}}
\newcommand{\course}{Complex Analysis}
\newcommand{\hwNum}{9}

\begin{document}

\pagestyle{fancy}
\fancyhead{}
\lhead{\zhstr{\name} (\uid)}
\chead{\course\ \#\hwNum}
\rhead{\today}

\section{Exercise 4.17}

Show that the sickle-shaped domain
$D = U_1(0) \setminus U_{1/2}(1/2)
= \{z \in \C : |z| < 1,\ |z - \tfrac{1}{2}| > \tfrac{1}{2}\}$
is an elementary domain.

\textit{Hint:} Consider the conformal map $z \mapsto \frac{1}{1-z}$.
What is the image of $D$ under this map?

\begin{solution}
  Consider the boundary of $D$. It contains two circles:
  \[
    \begin{aligned}
      &\{z = e^{i\theta} : \theta \in [0, 2\pi)\} \\
      &\{z = \frac{1}{2} + \frac{1}{2}e^{i\phi} : \phi \in [0, 2\pi)\}
    \end{aligned}
  \]
  We apply the conformal map $z \mapsto \frac{1}{1-z}$ to the
  first circle:
  \[
    \begin{aligned}
      &\left\{\frac{1}{1-e^{i\theta}} : \theta \in [0, 2\pi)\right\} \\
      =& \left\{\frac{1-e^{-i\theta}}{|1-e^{i\theta}|^2} : \theta \in
      [0, 2\pi)\right\} \\
      =& \left\{\frac{1-\cos\theta}{2(1-\cos\theta)} +
      i\frac{\sin\theta}{2(1-\cos\theta)} : \theta \in [0, 2\pi)\right\} \\
      =& \left\{\frac{1}{2} + i\frac{\sin\theta}{2(1-\cos\theta)} :
      \theta \in [0, 2\pi)\right\} \\
      =& \left\{\frac{1}{2} + i\frac{1}{2}\cot\frac{\theta}{2} :
      \theta \in [0, 2\pi)\right\}
    \end{aligned}
  \]
  which is a vertical line.  Similarly, we apply the conformal map
  to the second circle:
  \[
    \begin{aligned}
      &\left\{\frac{1}{1-\frac{1}{2} - \frac{1}{2}e^{i\phi}} : \phi
      \in [0, 2\pi)\right\} \\
      =& \left\{\frac{1}{\frac{1}{2}(1+e^{i\phi})} : \phi \in [0,
      2\pi)\right\} \\
      =& \left\{\frac{2}{1+e^{i\phi}} : \phi \in [0, 2\pi) \right\}
      % TODO: simplify the above expression.
    \end{aligned}
  \]
  which is the same vertical line.
\end{solution}

\section{Exercise 4.21}

Let $f$ be an analytic function in $U_R(0)$ with $R > 0$.
Show that if there exist constants $\rho \in (0,1)$ and $C > 0$
such that for sufficiently large $n \in \N$,
$|f(\tfrac{1}{n})| \le C\rho^n$,
then $f \equiv 0$.

\begin{solution}

\end{solution}

\section{Exercise 4.33}

Decide whether there are analytic functions
$f, g : U_R(0) \to \C$ with
\begin{enumerate}[(a)]
  \item $f(\tfrac{1}{n}) = f(-\tfrac{1}{n}) = \tfrac{1}{n^2}$,
    $n \ge 1$;
  \item $g(-\tfrac{1}{2n}) = g(\tfrac{1}{2n}) = \tfrac{1}{2n}$,
    $n \ge 1$.
\end{enumerate}

\begin{solution}

\end{solution}

\section{Exercise 4.34}

Let $f$ be an analytic function on a domain $D$ satisfying that for
any $z \in D$, there exists an integer $n = n(z) \ge 0$ such that
$f^{(n)}(z) = 0$.  Show that $f$ is a polynomial.

\textit{Hint:} $D$ can be represented as
$\bigcup_{n=0}^{\infty} g_n^{-1}(\{0\})$,
where $g_n(z) = f^{(n)}(z)$.

\begin{solution}

\end{solution}

\section{Exercise 4.36}

If $D$ is not connected, then $H(D)$ might have divisors of zero.
Give one such example.

\begin{solution}

\end{solution}

\section{Exercise 4.43}

Show that if $f : D \to D$ is analytic on a domain and satisfies
$f \circ f = f$, then $f$ is either constant or the identity map.

\begin{solution}

\end{solution}

\section{Exercise 4.46}

Let $w_1, \ldots, w_n$ be points on the unit circle.
Prove that there exists a point $z$ on the unit circle such that
the product of the distances from $z$ to the points $w_j$ is at
least~$1$, i.e.,
$\prod_{j=1}^{n} |z - w_j| \ge 1$.
Conclude that there exists a point $w_u$ on the circle such that
$\prod_{j=1}^{n} |w_u - w_j| = 1$.

\begin{solution}

\end{solution}

\section{Exercise 4.47}

Show that if the real part of an entire function $f$ is bounded,
then $f$ is constant.

\begin{solution}

\end{solution}

\section{Exercise 4.49}

Prove the fundamental theorem of algebra by Theorem~4.48
(Minimum modulus).

\begin{solution}

\end{solution}

\section{Exercise 4.50}

Let $f$ be nonconstant and analytic in an open set containing
$\overline{D}$.  Show that if $|f(z)| = 1$ for all $|z| = 1$,
then $D \subseteq f(D)$.

\textit{Hint:} One must show that $f(z) = w_0$ has a root for
every $w_0 \in D$.  Then it suffices to consider the case
$w_0 = 0$.  Why?

\begin{solution}

\end{solution}

\end{document}
